20 世纪 80 年代中期，格罗莫夫与埃利亚什伯格 (El87) 证明了著名的
$C^0$-辛刚性定理：\emph{任何非辛微分同胚}都不可能是一个辛微分同胚序列的\emph{（局部）一致极限}。这是\emph{辛拓扑}这一概念存在的第一个迹象。埃利亚什伯格
证明的主要思想依赖于构造某些在辛微分同胚下保持不变的
$C^0$-型不变量。随后，格罗莫夫在其奠基性论文 (Gr85) 中利用\emph{格罗莫夫非挤压定理}，完成了这种 $C^0$-型不变量的构造。埃克兰与
霍费尔 (EkH89)、(EkH90) 把这类 $C^0$-不变量形式化为\emph{辛容量}的概念。关于经典相空间
$\mathbb R^{2n}$ 上容量的完整论述，可参见 (HZ94)。霍费尔 (H90) 利用哈密顿流的几何，在任意辛流形上给出了一个非常典范的容量构造，即所谓的\emph{位移能}。这个构造建立在哈密顿微分同胚群
$\Ham(M,\omega)$ 上一个卓越的（无限维）双不变芬斯勒几何之上。相伴的双不变范数称为 $\Ham(M,\omega)$ 上的\emph{霍费尔范数}。霍费尔 (H90)、波尔捷罗维奇 (Po93) 以及弗朗索瓦·拉隆德（François Lalonde）与麦克达夫 (LM95a) 随后在逐渐增广的一般性下证明了该范数的非退化性。

本章假设 $M$ 或者是闭的，或者是开的；当 $M$ 非紧或有边界时，所有函数都紧支于其内部。

\section{哈密顿量的归一化}
\stfhSubIndex{I216}{gui^yi^hua}{归一化（normalization）}{ha^mi^dun^liang^de}{哈密顿量的（of Hamiltonian）}
\label{sec:normalization-hamiltonians}

我们先考虑一个不一定为辛流形的任意流形 $M$。

回忆一般李群理论：在任意李群 $G$ 上，李代数 $\mathfrak g$ 上任何\emph{在伴随作用下不变}的范数，都通过从单位元到 $G$ 的给定群元素的右平移，在切丛 $TG$ 上诱导一个双不变芬斯勒范数。

在 $\Diff(M)$ 的情形，李代数对应于向量场集合 $\mathfrak X(M)$，其李括号为
$[X,Y]=\mathcal L_X(Y)$。回忆 $\mathcal L_X(Y)$ 的定义是
\[
 \mathcal L_XY=\left.\frac d{dt}\right|_{t=0}(\phi_t)_*Y,
\]
其中 $\phi_t$ 是 $\dot x=X(x)$ 的流。$\Diff(M)$ 在 $\mathfrak X(M)$ 上的伴随作用由
\[
 \operatorname{Ad}_\phi(Y)
 =\left.\frac d{dt}\right|_{t=0}\phi\psi_t\phi^{-1}
\]
实现，其中 $\psi_t$ 是满足 $\psi_0=\operatorname{id}$ 且
$Y=\left.(d/dt)\right|_{t=0}\psi_t$ 的单参数子群。换言之，
$\operatorname{Ad}_\phi(Y)$ 是如下定义的向量场：
\begin{equation}
 \operatorname{Ad}_\phi(Y)(x)
 =d\phi\bigl(Y(\phi^{-1}(x)\bigr)
 =\phi_*(Y)(x),
 \label{eq:adjoint-action-diffeomorphism}
\end{equation}
它恰是 $Y$ 在 $\phi$ 下的推前。

现在考虑辛流形 $(M,\omega)$。通过考察 $\Ham(M,\omega)$ 中的光滑单参数子群，我们已经在第 2.3 节证明，$\Ham(M,\omega)$ 的李代数
$\operatorname{ham}(M,\omega)$ 关于由 $\omega$ 相伴的泊松括号典范同构于
$C^\infty(M)/\mathbb R$；更准确地说，赋值
\[
 f\longmapsto-X_f=X_{-f}
\]
是 $\bigl(C^\infty(M)/\mathbb R,\{\mathord\cdot,\mathord\cdot\}\bigr)$ 与
$\bigl(\mathfrak X(M),[\mathord\cdot,\mathord\cdot]\bigr)$ 之间的李代数同构。回忆恒等式
\[
 \psi^*(X_f)=X_{f\circ\psi}
\]
以及
\[
 \{f\circ\psi,g\circ\psi\}=\{f,g\}\circ\psi
\]
对任意辛微分同胚 $\psi$ 都成立。因此，对应于 $\psi^{-1}$ 在
$C_0^\infty(M)$ 上的伴随作用的拉回
$\psi^*X_f=(\psi^{-1})_*(X_f)$，恰由复合
\[
 f\circ\psi
\]
给出。显然，这一作用延拓到整个微分同胚群 $\Diff(M)$。

由于许多稍后将看到的原因，通过为 $\mathbb R$ 在 $C^\infty(M)$ 上加常数的作用选取一个良好截面，把商空间 $C^\infty(M)/\mathbb R$ 实现为
$C^\infty(M)$ 的子空间是有用的。这类截面有两个典范选择，它们在
$\Symp(M,\omega)$ 的伴随作用下不变。我们把每个截面中的任意元素称为\emph{归一化的}。现在详细说明这两种归一化，把讨论分为闭情形与开情形；后一种情形指 $M$ 非紧或具有非空边界。

先考虑 $(M,\omega)$ 闭的情形。令 $d\mu$ 是体积形式
$\omega^n/n!$ 在 $M$ 上诱导的刘维尔测度。依定义，该测度在
$\Symp(M,\omega)$ 下不变，并且在闭流形 $M$ 上有限。

\begin{definition}\label{def:mean-normalized-functions}
定义
\[
 C_0^\infty(M)
 =\left\{f\in C^\infty(M)\mathrel{\bigg|}
   \int_M f\,d\mu=0\right\}.
\]
我们称 $C_0^\infty(M)$ 中的任意元素为\emph{均值归一化的}。
\stfhSubIndex{I215}{gui^yi^hua}{归一化（normalization）}{jun^zhi^gui^yi^hua}{均值归一化（mean-normalized）}
\end{definition}

下面的结论是刘维尔引理，即引理 2.3.2 的一个简单推论。

\begin{lemma}\label{lem:poisson-bracket-zero-mean}
假设 $M$ 是闭的。则对所有 $f,g\in C^\infty(M)$，都有
\[
 \int_M\{f,g\}\,d\mu=0,
\]
从而
\[
 \{C^\infty(M),C^\infty(M)\}\subset C_0^\infty(M).
\]
特别地，$C_0^\infty(M)$ 在 $\Ham(M,\omega)$ 的伴随作用下不变。
\end{lemma}

\begin{proof}
利用 $\mathcal L_{X_g}\omega=0$ 与嘉当“魔术”公式，计算得
\[
\begin{aligned}
 \int_M\{f,g\}\,d\mu
 &=\int_M\mathcal L_{X_g}(f)\,d\mu
  =\int_M\mathcal L_{X_g}(f)\frac{\omega^n}{n!}\\
 &=\int_M\mathcal L_{X_g}\left(f\frac{\omega^n}{n!}\right)
  =\int_M d\left(X_g\mathbin{\lrcorner}
      \frac{f\omega^n}{n!}\right)=0.
\end{aligned}
\]
\end{proof}

下面考虑 $(M,\omega)$ 开的情形。文献中一直采用如下归一化。

\begin{definition}\label{def:compactly-supported-functions}
定义
\[
 C_c^\infty(M)
 =\{f\in C^\infty(M)\mid
   \operatorname{supp}f\text{ 是紧的并包含于 }\operatorname{Int}(M)\}.
\]
这里 $c$ 表示“紧支”。
\end{definition}

然而，从下面将详述的哈密顿几何观点看，这种归一化似乎并非最优选择。

在开流形 $(M,\omega)$ 上，自然应当考虑\emph{紧支}哈密顿微分同胚（更一般地，紧支辛微分同胚）。

\begin{definition}\label{def:weak-strong-compact-support}
\stfhIndex{I021}{00-symbol-021}{$\Ham^{sc}(M,\omega)$}
\stfhIndex{I022}{00-symbol-022}{$\Ham^{wc}(M,\omega)$}
\stfhIndex{I264}{qiang^jin^zhi}{强紧支（strongly compactly supported）}
\stfhIndex{I303}{ruo^jin^zhi}{弱紧支（weakly compactly supported）}
称紧支辛微分同胚 $\phi$ 为
\begin{enumerate}[label=(\arabic*)]
 \item \emph{弱紧支哈密顿微分同胚}，如果存在哈密顿量
 $H:[0,1]\times M\to\mathbb R$，使其哈密顿向量场 $X_H$ 紧支，且
 $\phi=\phi_H^1$；
 \item \emph{强紧支哈密顿微分同胚}，如果存在紧支哈密顿量
 $H:[0,1]\times M\to\mathbb R$，使得 $\phi=\phi_H^1$。
\end{enumerate}
相应的这两类微分同胚集合分别记为
$\Ham^{wc}(M,\omega)$ 与 $\Ham^{sc}(M,\omega)$。
\end{definition}

显然有包含关系
\[
 \Ham^{sc}(M,\omega)\subset\Ham^{wc}(M,\omega)
 \subset\Symp^c(M,\omega).
\]
现在分析这些李群各自的李代数。显然，
$\Symp_0^c(M,\omega)$ 的李代数是紧支局部哈密顿向量场的集合。其次，依定义有
\[
 \Ham^{sc}(M,\omega)
 =\{\phi\in\Symp^c(M,\omega)\mid
   \phi=\phi_H^1,\ \operatorname{supp}(H)\subset\operatorname{Int}(M)
   \text{ 是紧的}\},
\]
其中 $\operatorname{supp}(H)$ 定义为
\stfhIndex{I044}{00-symbol-044}{$\operatorname{supp}(H)$}
\[
 \operatorname{supp}(H)=\bigcup_{t\in[0,1]}\operatorname{supp}(H_t).
\]
容易看出，$\Ham^{sc}(M,\omega)$ 的李代数由
\[
 \operatorname{ham}^{sc}(M,\omega)=C_c^\infty(M)
\]
给出。

\stfhIndex{I114}{duan}{端（ends）}
最后考虑 $\Ham^{wc}(M,\omega)$。这需要讨论 $M$ 的端结构。回忆，如果存在递增的紧子集序列
\begin{equation}
 M_1\subseteq M_2\subseteq\cdots\subseteq M_i\subseteq\cdots
 \label{eq:exhaustion-ends}
\end{equation}
使得 $\bigcup_iM_i=M$，并且存在 $N_0\in\mathbb Z_+$，使得
$M\setminus M_i$ 的连通分支数对所有 $i\geq N_0$ 保持不变，就称
$M$ 具有\emph{有限多个端}。我们把这个常数称为 $M$ 的\emph{端数}。容易检验，这个数与穷竭序列\eqref{eq:exhaustion-ends}的选择无关。以
$\infty_\alpha$ 表示 $M$ 的一个端。

\begin{definition}\label{def:average-constant-at-ends}
\stfhIndex{I007}{00-symbol-007}{$C_{ac}^{\infty}(M)$}
假设 $M$ 具有有限多个端。定义
\[
 C_{ac}^\infty(M)
 =\left\{f\in C^\infty(M)\mathrel{\bigg|}
   \sum_{\alpha\in\operatorname{End}(M)}f(\infty_\alpha)=0\right\}.
\]
\end{definition}

容易检验，对于任意紧支哈密顿向量场 $X$，总可以给它的生成哈密顿量加上一个常数，使新的哈密顿量属于 $C_{ac}^\infty(M)$。因此，
$C_{ac}^\infty(M)$ 典范同构于 $\Ham^{wc}(M,\omega)$ 的李代数。于是
$C_{ac}^\infty(M)$ 在 $\Ham^{wc}(M)$ 下不变，并且
\begin{equation}
 \dim\bigl(C_{ac}^\infty(M)/C_c^\infty(M)\bigr)
 =\#\operatorname{End}(M)-1.
 \label{eq:ends-quotient-dimension}
\end{equation}

下面的例子表明
$\Ham^{sc}(M,\omega)\subsetneq\Ham^{wc}(M,\omega)$，即存在并非强紧支的弱紧支哈密顿微分同胚。由\eqref{eq:ends-quotient-dimension}可知，要存在这样的微分同胚，$M$ 的端数必须至少为 $2$。

\begin{example}\label{ex:cylinder-weak-not-strong}
考虑柱面 $M=\mathbb R\times S^1$ 及辛形式
$\omega=d\tau\wedge dt$。注意 $\mathbb R\times S^1$ 的端数为 $2$。考虑由
\[
 f(\tau,t)=\rho(\tau)
\]
定义的哈密顿量 $f:\mathbb R\times S^1\to\mathbb R$，其中
$\rho:\mathbb R\times[0,1]$ 是满足
\begin{equation}
 \rho(\tau)=
 \begin{cases}
  -\dfrac12,&\tau\leq-K,\\[0.3em]
   \dfrac12,&\tau\geq K
 \end{cases}
 \label{eq:cylinder-end-profile}
\end{equation}
的单调函数，这里 $K>0$。相应的哈密顿向量场 $X_\rho$ 紧支，且其流
$\phi_\rho^t$ 属于 $\Ham^{wc}(M,\omega)$。
\end{example}

\begin{exercise}\label{ex:cylinder-flow-not-strong}
证明当 $t\ne0$ 时，$\phi_\rho^t$ 不属于 $\Ham^{sc}(M,\omega)$。
\end{exercise}
